Using Nauty, 
Orbiter can compute canonical forms and 
automorphism groups of codes.
For linear codes, 
the semilinear automorphism group can be computed.

\bigskip


Consider the $[3,2,2]$ code generated by 
$$
\left[
\begin{array}{ccc}
	1&0&1\\
	0&1&1\\
\end{array}
\right]
$$
Using the following command, 
the semilinear automorphism group of the code 
can be computed:

\sourcecode{code_3_2_aut}



The code has a semilinear automorphism group of order 6.
The following report is written:
\orbiteroutput{GRAPHICS/LATEX/summary_of_orbits_code_3_2}

\bigskip




The command 

\sourcecodevariable{CODE_RM_3_1_GENMA}


\sourcecode{RM_3_1_group}

computes the automorphism group of the Reed-Muller code. 
The group is an affine linear group $\AGL(3,2)$ of order 1344. 
A report is created, showing the automorphism group 
and the action on $\PG(3,2),$ with the Reed-Muller code distinguished.


\bigskip

The following command creates a drawing of the incidence matrix between points 
and lines in $\PG(3,2),$ 
with the Reed-Muller code distinguished:
\sourcecode{RM_3_1_group_and_diagram}

The drawing is shown below:
\orbiteroutput{GRAPHICS/WRAPPERS/RM_3_1_object0_TDA_flag_orbits_draw}


\bigskip




The command 

\sourcecode{RS_6_4_7_group}

shows that the automorphism group has order $12$.
After some shortening, the output is:
\orbiteroutput{GRAPHICS/LATEX/RS_6_4_7}

\bigskip


The command 

\sourcecode{GV_n15_k6_d5_group}

computes the automorphism group of the Gilbert-Varshamov code from Section~\ref{sec:Bounds}. 
It has order $12$.

\bigskip

